\documentclass[12pt,a4paper]{article}
\usepackage{mathwizards-ingles}
\title{%
  \vspace{-1cm}
  {\Huge\bfseries\textcolor{mes2}{Mes 2 · Semana 6}}\\[0.3cm]
  {\Large Present Continuous, Descripciones y Minería Activa}\\[0.2cm]
  {\large\color{grissuave} Nivel A1 $\rightarrow$ A2 · Inicio de la Minería de Oraciones}\\[0.5cm]
  \rule{\textwidth}{1.5pt}
}
\author{}
\date{}

\begin{document}
\maketitle
\thispagestyle{empty}

\begin{cajanivel}
\textbf{Objetivo:} El estudiante inicia la \textbf{minería activa de oraciones}, extrayendo frases del input que consume. Se introduce el \textbf{Present Continuous} y su contraste con el Present Simple. Se expanden las descripciones con \textbf{adjetivos} y se incorporan chunks de \textbf{opinión y supervivencia} para viajes/compras.
\end{cajanivel}

% ═══════════════════════════════════════════════════════════════════════════════
\section{Gramática Emergente: Present Continuous}
% ═══════════════════════════════════════════════════════════════════════════════

\begin{cajagrammar}[Present Continuous: \textit{to be} + verbo + \textbf{-ing}]
Se usa para hablar de acciones que están \textbf{ocurriendo ahora mismo} o \textbf{temporalmente}.

\medskip
\begin{tabularx}{\textwidth}{@{} l X @{}}
\toprule
\textbf{Estructura} & \textbf{Ejemplo} \\
\midrule
\textbf{Afirmativo:} Sujeto + am/is/are + V-ing & \textit{She \textbf{is reading} a book right now.} \\
\textbf{Negativo:} Sujeto + am/is/are + not + V-ing & \textit{I'\textbf{m not watching} TV.} \\
\textbf{Pregunta:} Am/Is/Are + sujeto + V-ing? & \textit{\textbf{Are} you \textbf{studying} English?} \\
\bottomrule
\end{tabularx}

\medskip
\begin{cajaejemplo}[Formación del -ing]
\begin{tabularx}{\textwidth}{@{} l l l @{}}
\textbf{Regla} & \textbf{Base} & \textbf{-ing} \\
\midrule
General: + ing & \textit{eat} & \textit{eat\textbf{ing}} \\
Termina en -e muda: quitar e + ing & \textit{make} & \textit{mak\textbf{ing}} \\
CVC corta: doblar consonante + ing & \textit{run} & \textit{runn\textbf{ing}} \\
Termina en -ie: ie $\rightarrow$ y + ing & \textit{lie} & \textit{ly\textbf{ing}} \\
\end{tabularx}
\end{cajaejemplo}
\end{cajagrammar}

\begin{cajagrammar}[Contraste: Present Simple vs. Present Continuous]
Este contraste es una fuente constante de confusión. Hay que adquirirlo por \textbf{exposición}, no por reglas.

\medskip
\begin{tabularx}{\textwidth}{@{} l X X @{}}
\toprule
& \textbf{Present Simple} & \textbf{Present Continuous} \\
\midrule
\textbf{Uso} & Rutinas, hábitos, hechos & Ahora mismo, temporal \\
\textbf{Señales} & always, usually, every day & now, right now, at the moment \\
\textbf{Ejemplo} & \textit{I \textbf{drink} coffee every morning.} & \textit{I'\textbf{m drinking} coffee right now.} \\
\bottomrule
\end{tabularx}

\medskip
\begin{cajaejemplo}[Verbos que NO se usan en Continuous (Stative Verbs)]
Algunos verbos describen estados, no acciones, y \textbf{no llevan -ing}:
\begin{multicols}{3}
\begin{itemize}[nosep, leftmargin=0.5cm]
  \item \eng{know} \esp{saber}
  \item \eng{want} \esp{querer}
  \item \eng{need} \esp{necesitar}
  \item \eng{like} \esp{gustar}
  \item \eng{love} \esp{amar}
  \item \eng{hate} \esp{odiar}
  \item \eng{believe} \esp{creer}
  \item \eng{understand} \esp{entender}
  \item \eng{remember} \esp{recordar}
\end{itemize}
\end{multicols}
$\checkmark$ \textit{I \textbf{know} the answer.} \quad $\times$ \textit{I'm knowing the answer.}
\end{cajaejemplo}
\end{cajagrammar}

% ═══════════════════════════════════════════════════════════════════════════════
\section{Chunks: Adjetivos y Orden de Adjetivos}
% ═══════════════════════════════════════════════════════════════════════════════

\begin{cajachunk}[Common Adjectives — Adjetivos de Alta Frecuencia]
\begin{multicols}{3}
\begin{tabularx}{\linewidth}{@{} l l @{}}
\eng{big} & \esp{grande} \\
\eng{small} & \esp{pequeño} \\
\eng{old} & \esp{viejo} \\
\eng{new} & \esp{nuevo} \\
\eng{good} & \esp{bueno} \\
\eng{bad} & \esp{malo} \\
\eng{happy} & \esp{feliz} \\
\eng{sad} & \esp{triste} \\
\end{tabularx}

\begin{tabularx}{\linewidth}{@{} l l @{}}
\eng{hot} & \esp{caliente} \\
\eng{cold} & \esp{frío} \\
\eng{fast} & \esp{rápido} \\
\eng{slow} & \esp{lento} \\
\eng{easy} & \esp{fácil} \\
\eng{hard} & \esp{difícil} \\
\eng{beautiful} & \esp{bonito} \\
\eng{ugly} & \esp{feo} \\
\end{tabularx}

\begin{tabularx}{\linewidth}{@{} l l @{}}
\eng{long} & \esp{largo} \\
\eng{short} & \esp{corto} \\
\eng{tall} & \esp{alto} \\
\eng{cheap} & \esp{barato} \\
\eng{expensive} & \esp{caro} \\
\eng{interesting} & \esp{interesante} \\
\eng{boring} & \esp{aburrido} \\
\eng{important} & \esp{importante} \\
\end{tabularx}
\end{multicols}
\end{cajachunk}

\begin{cajagrammar}[Comparatives \& Superlatives — Introducción]
\begin{tabularx}{\textwidth}{@{} l l l l @{}}
\toprule
\textbf{Regla} & \textbf{Adjetivo} & \textbf{Comparativo} & \textbf{Superlativo} \\
\midrule
1 sílaba: +er / +est & \textit{tall} & \textit{tall\textbf{er}} & \textit{the tall\textbf{est}} \\
1 sílaba CVC: doblar +er & \textit{big} & \textit{bigg\textbf{er}} & \textit{the bigg\textbf{est}} \\
Termina en -y: y $\rightarrow$ ier & \textit{happy} & \textit{happ\textbf{ier}} & \textit{the happ\textbf{iest}} \\
2+ sílabas: more / most & \textit{beautiful} & \textit{\textbf{more} beautiful} & \textit{the \textbf{most} beautiful} \\
\midrule
\multicolumn{4}{@{}l}{\textbf{Irregulares:}} \\
& \textit{good} & \textit{\textbf{better}} & \textit{the \textbf{best}} \\
& \textit{bad} & \textit{\textbf{worse}} & \textit{the \textbf{worst}} \\
\bottomrule
\end{tabularx}

\medskip
\begin{cajaejemplo}[Estructuras Clave]
\textbf{Comparativo:} A + is + \textbf{adj-er / more adj} + \textbf{than} + B\\
\textit{My house is \textbf{bigger than} your house.}

\medskip
\textbf{Superlativo:} A + is + \textbf{the adj-est / the most adj}\\
\textit{She is \textbf{the tallest} person in the class.}
\end{cajaejemplo}
\end{cajagrammar}

% ═══════════════════════════════════════════════════════════════════════════════
\section{Chunks: Opiniones}
% ═══════════════════════════════════════════════════════════════════════════════

\begin{cajachunk}[Opinion Chunks — Expresar Opiniones Simples]
\begin{tabularx}{\textwidth}{@{} X l @{}}
\eng{I think...} & \esp{Creo que...} \\
\eng{In my opinion, ...} & \esp{En mi opinión, ...} \\
\eng{I like ... because ...} & \esp{Me gusta ... porque ...} \\
\eng{I don't like ... because ...} & \esp{No me gusta ... porque ...} \\
\eng{I prefer ... to ...} & \esp{Prefiero ... a ...} \\
\eng{I agree.} & \esp{Estoy de acuerdo.} \\
\eng{I don't agree.} / \eng{I disagree.} & \esp{No estoy de acuerdo.} \\
\eng{That's a good idea.} & \esp{Es una buena idea.} \\
\eng{I'm not sure.} & \esp{No estoy seguro.} \\
\eng{Maybe.} / \eng{Perhaps.} & \esp{Quizás. / Tal vez.} \\
\end{tabularx}
\end{cajachunk}

% ═══════════════════════════════════════════════════════════════════════════════
\section{Chunks: Supervivencia — Compras y Viajes}
% ═══════════════════════════════════════════════════════════════════════════════

\begin{cajachunk}[Shopping — Compras]
\begin{tabularx}{\textwidth}{@{} X l @{}}
\eng{How much is this?} / \eng{How much does this cost?} & \esp{¿Cuánto cuesta esto?} \\
\eng{I'd like ... , please.} & \esp{Quisiera ..., por favor.} \\
\eng{Can I have ... ?} & \esp{¿Me puede dar ...?} \\
\eng{Do you have ... ?} & \esp{¿Tiene ...?} \\
\eng{I'm just looking, thanks.} & \esp{Solo estoy mirando, gracias.} \\
\eng{I'll take it.} & \esp{Me lo llevo.} \\
\eng{Can I pay by card?} & \esp{¿Puedo pagar con tarjeta?} \\
\eng{Here you go.} & \esp{Aquí tiene.} \\
\eng{Keep the change.} & \esp{Quédese con el cambio.} \\
\end{tabularx}
\end{cajachunk}

\begin{cajachunk}[Travel — Viajes]
\begin{tabularx}{\textwidth}{@{} X l @{}}
\eng{Where is the ... ?} & \esp{¿Dónde está el/la ...?} \\
\eng{How do I get to ... ?} & \esp{¿Cómo llego a ...?} \\
\eng{Is it far from here?} & \esp{¿Está lejos de aquí?} \\
\eng{Turn left / right.} & \esp{Gira a la izquierda / derecha.} \\
\eng{Go straight ahead.} & \esp{Sigue derecho.} \\
\eng{It's next to the ...} & \esp{Está al lado del/de la ...} \\
\eng{I'm lost.} & \esp{Estoy perdido/a.} \\
\eng{Can you help me?} & \esp{¿Puedes ayudarme?} \\
\eng{I need a taxi, please.} & \esp{Necesito un taxi, por favor.} \\
\end{tabularx}
\end{cajachunk}

% ═══════════════════════════════════════════════════════════════════════════════
\section{Oraciones Listas para Minar}
% ═══════════════════════════════════════════════════════════════════════════════

\begin{cajamineria}[Oraciones \textit{i+1} — Semana 6]

\begin{enumerate}[leftmargin=1.2cm]
  \item \textit{She \underline{is reading} a new book right now.}\\
    {\small\textbf{Objetivo:} reconocer la estructura del Present Continuous (is + V-ing).}

  \item \textit{I usually eat at home, but today I'\underline{m eating} at a restaurant.}\\
    {\small\textbf{Objetivo:} contraste Present Simple (usually eat) vs. Continuous (today I'm eating).}

  \item \textit{This shirt is \underline{cheaper} than that one.}\\
    {\small\textbf{Objetivo:} la estructura comparativa ``cheaper than''.}

  \item \textit{She is \underline{the tallest} girl in her class.}\\
    {\small\textbf{Objetivo:} la estructura superlativa ``the tallest''.}

  \item \textit{I \underline{think} this movie is very interesting.}\\
    {\small\textbf{Objetivo:} el chunk de opinión ``I think'' + stative verb.}

  \item \textit{How much \underline{does} this cost?}\\
    {\small\textbf{Objetivo:} la estructura interrogativa de compras con ``does''.}

  \item \textit{Excuse me, how do I \underline{get to} the train station?}\\
    {\small\textbf{Objetivo:} la expresión ``get to'' (llegar a).}

  \item \textit{He is \underline{running} in the park because it's a beautiful day.}\\
    {\small\textbf{Objetivo:} formación del -ing con duplicación de consonante (run $\rightarrow$ running).}

  \item \textit{I \underline{prefer} coffee to tea.}\\
    {\small\textbf{Objetivo:} la estructura ``prefer X to Y''.}

  \item \textit{Are you \underline{looking for} something?}\\
    {\small\textbf{Objetivo:} el phrasal verb ``look for'' (buscar) en continuous.}
\end{enumerate}
\end{cajamineria}

\vfill
\begin{center}
\rule{0.5\textwidth}{0.5pt}\\[0.3cm]
{\small\color{grissuave} Curso de Inmersión en Inglés · Mes 2 · Semana 6 · Nivel A1 $\rightarrow$ A2}
\end{center}

\end{document}
